\BourbakiExerciseGroup

\begin{enumerate}
\def\labelenumi{\arabic{enumi})}
\item
  Suppose A is an integral domain; let \(f\) be a polynomial of the ring \(\mathbf{A}[X_1, X_2, \ldots, X_n]\) of degree \(\leqslant k_i\) in \(X_i\) (for \(1 \leqslant i \leqslant n\) ). For each value of \(i\) ( \(1 \leqslant i \leqslant n\) ) let \(H_i\) be a set of \(k_i + 1\) elements of A. Show that if \(f(x_1, x_2, \ldots, x_n) = 0\) for each element \((x_i) \in \prod_{i=1}^{n} H_i\) , then \(f = 0\) .
\item
  Assume that A is an infinite integral domain; let \(\Phi\) be a set of polynomials \(\neq0\) in the ring \(A[X_{1},X_{2},\ldots,X_{n}]\) . Show that if the power of \(\Phi\) is strictly less than that of A, then there exists a subset H of \(A^{n}\) equipotent with A such that for each \(x=(x_{i})\in H\) and for each \(f\in\Phi\) we have \(f(x)\neq0\) .
\item
  Suppose that A is an integral domain and let B be an infinite subset of A. Show that if the polynomial \(f \in A[X]\) has degree \textgreater0, then the image of B under the polynomial function \(x \mapsto f(x)\) is equipotent with B.
\item
  Suppose that the additive group of A contains an infinite subgroup G whose non-zero elements are not divisors of zero in A. Show that the mapping \(f \mapsto \tilde{f}\) of \(A{[}X_{1}, \ldots, X_{p}{]}\) into the algebra of mappings of \(A^{p}\) into A is injective (observe that a polynomial of degree n in one indeterminate cannot have more than n distinct roots lying in G). This holds in particular when there is in A an element \(x_{0}\) , not a divisor of 0 and of infinite order in the additive group of A.
\item
  Let K be a field of infinitely many elements and let Q be the quaternion algebra over K, of type \((-1,0,-1)\) (III, p.~445). Show that the polynomial \(X^{2}+1\) has infinitely many roots in Q.
\item
  Let K be a finite field with q elements.
\end{enumerate}

\begin{enumerate}
\def\labelenumi{\alph{enumi})}
\tightlist
\item
  In the ring \(\mathbf{K}[\mathbf{X}_1, \mathbf{X}_2, \ldots, \mathbf{X}_n]\) let \(\mathfrak{a}\) be the ideal generated by the \(n\) polynomials
\end{enumerate}

\(X_{i}^{q}-X_{i}\quad(1\leqslant i\leqslant n)\) . Show that if \(f\in\mathfrak{a}\) , then \(f(x_{1},x_{2},\ldots,x_{n})=0\) for all \((x_{i})\in\mathbf{K}^{n}\) (observe that the multiplicative group \(K^{*}\) is of order q-1).

\begin{enumerate}
\def\labelenumi{\alph{enumi})}
\setcounter{enumi}{1}
\item
  If \(f\) is any polynomial of \(\mathbf{K}[\mathbf{X}_1, \mathbf{X}_2, \ldots, \mathbf{X}_n]\) , show that there exists precisely one polynomial \(\vec{f}\) such that \(\deg_i \vec{f} \leqslant q - 1\) for \(1 \leqslant i \leqslant n\) and such that \(f \equiv \vec{f} \pmod{\mathfrak{a}}\) . We thus have \(\deg \vec{f} \leqslant \deg f\) ; if \(f\) is such that \(f(x_1, x_2, \ldots, x_n) = 0\) for each \((x_i) \in \mathbf{K}^n\) , \(f\) belongs to the ideal \(\mathfrak{a}\) , which is thus the inverse image of 0 under the homomorphism \(f \mapsto \vec{f}\) (use Exercise 1).
\item
  Let \(f_1, f_2, \ldots, f_m\) be non-zero polynomials of \(\mathbf{K}[\mathbf{X}_1, \mathbf{X}_2, \ldots, \mathbf{X}_n]\) such that \(f_i(0, 0, \ldots, 0) = 0 (1 \leqslant i \leqslant m)\) and the sum of the total degrees of the \(f_i\) is \(< n\) . Show that there exists an element \((x_1, x_2, \ldots, x_n) \in \mathbf{K}^n\) distinct from \((0, 0, \ldots, 0)\) such that
\end{enumerate}

\[
f _ {i} (x _ {1}, x _ {2}, \dots , x _ {n}) = 0 \quad \text { for } \quad 1 \leqslant i \leqslant m.
\]

(Observe that if this were not so, the polynomial \(\prod_{i=1}^{m} (1 - f_i^{q-1})\) would be congruent mod a to the polynomial \(\prod_{i=1}^{n} (1 - X_i^{q-1})\) , and use \(b\) ).

\begin{enumerate}
\def\labelenumi{\arabic{enumi})}
\setcounter{enumi}{6}
\item
  Let K be a finite field with q elements, and B the product ring \(K^{1}\) where I is an infinite set. Give an example of a non-zero polynomial f of B{[}X{]} such that \(f(x)=0\) for all \(x\in B\) .
\item
  Let B be the additive group \(\mathbf{Z} \oplus (\mathbf{Z}/2\mathbf{Z})\) , with the commutative ring structure obtained by putting \((a, b) \cdot (a', b') = (aa', ab' + ba')\) for \((a, b) \in \mathbf{B}\) and \((a', b') \in \mathbf{B}\) . Let \(\varepsilon = (0, 1) \in \mathbf{B}\) , and \(f(\mathbf{X}) = \varepsilon \mathbf{X}(\mathbf{X} - 1) \in \mathbf{B}[\mathbf{X}]\) . Then \(f(\alpha) = 0\) for all \(\alpha \in B\) .
\item
  Let \(k\) be an integer \(> 0\) such that \(k \cdot 1 = 0\) in \(\mathbf{A}\) . Let \(h\) be an integer \(> 0\) and \(f\) the polynomial \((\mathbf{X} - \alpha)^{k + h} + (\mathbf{X} - \alpha)^k\) . Then \(\alpha\) is a root of order \(k\) of \(f\) and a root of order \(\geqslant k + h - 1\) of \(Df\) .
\end{enumerate}

10) Let \(f\) be a non-zero element of \(\mathbf{Z}[\mathbf{X}_1, \ldots, \mathbf{X}_n]\) . Let \(\mathrm{N}(q)\) be the number of zeros of \(f\) in the cell \(|x_1| \leqslant q, \ldots, |x_n| \leqslant q\) . Then \(\mathrm{N}(q) / q^n\) tends to 0 as \(q \to +\infty\) .

\begin{enumerate}
\def\labelenumi{\arabic{enumi})}
\setcounter{enumi}{10}
\item
  Let P be the set of all prime numbers, \(P'\) an infinite subset of P and \(f \in \mathbf{Z}[(\mathbf{X}_i)]\) . For each \(p \in P'\) let \(f_p\) be the element of \((\mathbf{Z}/p\mathbf{Z})[(\mathbf{X}_i)]\) obtained by applying to the coefficients of f the canonical homomorphism of Z onto Z/pZ. Suppose that for each \(p \in P'\) and each \(x \in (\mathbf{Z}/p\mathbf{Z})^l\) we have \(f_p(x) = 0\) ; then f = 0.
\item
  Let \(\mathbf{M}\) be an A-module, and let \(\alpha_0, \ldots, \alpha_n \in \mathbf{A}\) be such that for \(i \neq j\) the homothety with respect to \(\alpha_i - \alpha_j\) in \(\mathbf{M}\) is bijective. For any \(m_0, \ldots, m_n \in M\) there exists one and only one \(f \in \mathbf{M}[\mathbf{X}]\) such that \(f(\alpha_i) = m_i\) for \(0 \leqslant i \leqslant n\) and \(\deg f \leqslant n\) .
\item
  Let \(\alpha_{i}\) ( \(1 \leqslant i \leqslant n\) ) be \(n\) distinct elements of a commutative field \(\mathbf{K}\) , \(\beta_{i}\) ( \(1 \leqslant i \leqslant n\) ) and \(\gamma_{i}\) ( \(1 \leqslant i \leqslant n\) ) \(2n\) arbitrary elements of \(\mathbf{K}\) . Show that there exists one and only one polynomial \(f \in \mathbf{K}[\mathbf{X}]\) of degree \(\leqslant 2n - 1\) such that \(f(\alpha_{i}) = \beta_{i}\) and \(f'(\alpha_{i}) = \gamma_{i}\) for \(1 \leqslant i \leqslant n\) (Hermite interpolation formula) (to begin with consider the case where \(2n - 1\) of the \(2n\) elements \(\beta_{i}, \gamma_{i}\) are zero).
\end{enumerate}

